% SGA 6, Exposición XIII: fragmento español alineado con la fuente.
% Alcance exacto: fr.tex líneas 15151--15300. El final deja abierta la
% enumeración del Teorema 6.7 para continuarla en el fragmento siguiente.
% Las correcciones matemáticamente forzadas se conservan mediante \sisi.

\providecommand{\sisi}[2]{#2}

\begin{enumerate}[label=(\alph*)]
\item Toda generalización de $s$ permanece en $S_0$.
\item Toda especialización de $s$, próxima a $s$, permanece en $S_0$.
\item Bajo las condiciones de (ii), toda especialización de $s$ permanece
en $S_0$.
\end{enumerate}

Sea $f:X'\to X$ una presentación de Chow: $X'$ es proyectivo sobre $S$ y
$f$ es sobreyectivo. Si se admiten (a), (b), (c) y (iii) para $X'$, entonces
resultan para $X$ en virtud de 4.3 y 3.5. Se puede suponer, por tanto, que
$X$ es proyectivo sobre $S$.

Sea $H$ amplio para $X/S$. Si, para un $s\in S$,
$(L^{\otimes m}\otimes H)_s$ es amplio, entonces
$(L^{\otimes m}\otimes H)_t$ es amplio para toda generalización $t$ de $s$
(EGA III 4.7.1). Gracias a la equivalencia de 4.6 (i) y (iv), queda así
establecido (a).

Para (b), (c) y 4.7 (iii), se puede suponer que $S$ es íntegro, con $s$
como punto genérico en (b) y (c). Para (b) y (c) esto es claro, y para (iii)
resulta de 3.4 (i) aplicado al morfismo $\coprod S_i\to S$, donde los $S_i$
son las componentes irreducibles de $S$ dotadas de la estructura reducida
inducida.

Existen un esquema íntegro $S'$ y un morfismo f.p.p.f. $S'\to S$ tales que
toda componente irreducible $X'_i$ de $X_{S'}$, dotada de la estructura
reducida inducida, es íntegra sobre $S'$ (XII 1.3). Para demostrar (b) y
(iii), basta hacerlo para $X_{S'}/S'$: en (b), por descenso; en (iii), por
recurrencia noetheriana y 3.4 (i).

Consideremos finalmente el morfismo sobreyectivo canónico
$\coprod X'_i\to X_{S'}$. Entonces (b) y (iii) para $X_{S'}/S'$ resultan de
(c) y (iii) para los $X'_i/S'$, en virtud de 4.3, 3.5 y 3.7.

La cuestión se ha reducido, pues, a demostrar (c) y (iii), suponiendo que
$S$ es íntegro de punto genérico $s$, y que $X$ es proyectivo e íntegro
sobre $S$.

Sea $\mathcal O_X(1)$ un haz muy amplio para $X/S$. Para todos $m,n$,
$\chi(L_t^{\otimes m}(n))$ y $\chi(\mathcal O_{X_t}(n))$ no dependen de
$t\in S$ (EGA III 7.9.4). Por tanto, (c) resulta de
4.6 (i) $\Longleftrightarrow$ (v), que implica también que la familia de las
clases de haces invertibles $\tau$-equivalentes a $0$ sobre las fibras de
$X/S$ tiene un único polinomio de Hilbert; (iii) resulta, pues, de 2.11.

\textbf{Observación 4.8.} Sea $X$ propio e íntegro sobre un esquema
noetheriano $S$. Sin hipótesis de proyectividad sobre $X/S$, se ignora si
sigue siendo cierto que
$\underline{\Pic}^\tau_{X/S}\to\underline{\Pic}_{X/S}$ es una inmersión
\emph{cerrada}, salvo cuando $X$ es normal sobre $S$ y
$\underline{\Pic}_{X/S}$ es \emph{representable} (FGA 236, 2.3).

% Añadido en las páginas 654--669
\providecommand{\NS}{\operatorname{NS}}
\providecommand{\Picard}{\operatorname{Pic}}
\providecommand{\Et}{\mathrm{et}}
\providecommand{\rk}{\operatorname{rk}}
\providecommand{\im}{\operatorname{Im}}
\providecommand{\HH}{\mathrm H}

\subsection*{5. Teoremas de finitud de tipo «Néron--Severi»}

\textbf{Teorema 5.1.} \emph{Sea $X$ propio sobre un esquema noetheriano
$S$. Entonces los grupos de Néron--Severi
$\underline{\Pic}_{X_K/K}(K)/\underline{\Pic}^0_{X_K/K}(K)$ de las fibras
geométricas $X_K/K$ de $X/S$ son de tipo finito, y su rango y el orden de su
subgrupo de torsión
$\underline{\Pic}^\tau_{X_K/K}(K)/\underline{\Pic}^0_{X_K/K}(K)$
permanecen acotados superiormente.}

En efecto, por recurrencia noetheriana basta demostrar 5.1 después de
reemplazar $S$ por un abierto no vacío arbitrario de $S_{\red}$; así se
puede suponer que $\underline{\Pic}^\tau_{X/S}$ está representado por un
$S$-esquema de tipo finito (XII 1.2 y 4.7 (iii)). Entonces el orden de
$\underline{\Pic}^\tau_{X_K/K}(K)/\underline{\Pic}^0_{X_K/K}(K)$ es igual
al número de componentes conexas de
$(\underline{\Pic}^\tau_{X/S})_K$, número que permanece constante sobre un
abierto no vacío de $S$ (EGA IV 9.7.9).

Queda acotar superiormente el rango de
$\underline{\Pic}_{X_K/K}(K)/\underline{\Pic}^\tau_{X_K/K}(K)$. Como en la
demostración de 4.7 (iii), la cuestión se reduce fácilmente al caso en que
$S$ es íntegro y $X$ es proyectivo e íntegro sobre $S$. Sea \(r\) la dimensión
común de las fibras de $X/S$. Si \(r\) es \(0\) (resp. \(1\)), este rango es
evidentemente $0$ (resp. $1$). Si \(r\) es \(\geq2\), se reducirá la cuestión al caso
$r=2$ y $X$ liso sobre $S$. Tras restringir $S$, se puede encontrar una
sección lineal $Y$ de $X$ que sigue siendo íntegra sobre $S$, con fibras de
dimensión $2$; en virtud de 4.3, se puede reemplazar $X$ por $Y$. Finalmente,
en virtud de 4.3 y 3.10, utilizando la resolución de singularidades de las
superficies, y tras reemplazar $S$ por un recubrimiento plano de un abierto
no vacío de $S$, se puede reemplazar $X$ por un $S$-esquema liso.

Ahora, gracias a 4.6 (i) $\Longleftrightarrow$ (iii) y al teorema de finitud
y especialización de los grupos de cohomología étale (SGA 4 XVI 2.2), 5.1
resulta de la proposición siguiente.

\textbf{Proposición 5.2.} \emph{Sean $k$ un cuerpo algebraicamente cerrado y
$X$ un $k$-esquema propio y liso. Entonces el grupo $C^p_{\rm num}(X)$ de
los ciclos algebraicos de codimensión $p$ sobre $X$, módulo equivalencia
numérica, es libre de rango $\leq b^{2p}$, el número de Betti.}

En efecto, si se omite la acción de las raíces de 1, la cohomología étale
$H^*(X,\Q_\ell)=_{\rm dfn}
\left[\varprojlim_\nu H^*(X,\Z/\ell^\nu\Z)\right]
\otimes_{\Z_\ell}\Q_\ell$, donde $(\ell,\operatorname{car}(k))=1$, se
convierte en una cohomología de Weil: satisface la dualidad de Poincaré y la
fórmula de Künneth, y recibe un homomorfismo funtorial no nulo del anillo de
ciclos algebraicos módulo equivalencia racional $C^*_{\rm rat}(X)$ (véase X
7.13).

Sean $a_1,\ldots,a_m\in C^{r-p}_{\rm rat}(X)$, donde $r=\dim(X)$, elementos
cuyas imágenes forman una base del $\Q_\ell$-espacio vectorial en
$H^{2(r-p)}(X,\Q_\ell)$ engendrado por los ciclos algebraicos. Consideremos
el homomorfismo
\[
\alpha:C^p_{\rm rat}(X)\longrightarrow \Z^m
\]
definido por
$\alpha(b)=(\langle b.a_1\rangle,\ldots,
\langle b\sisi{,}{.}a_m\rangle)$, donde $\langle\ \rangle$ denota el número
de intersección. Se ve fácilmente que la imagen de $\alpha$ es
$C^p_{\rm num}(X)$.

\textbf{Observación 5.3.} Sean $k$ un cuerpo algebraicamente cerrado y $X$
una superficie propia y lisa sobre $k$. El segundo número de Betti $b^2(X)$
puede interpretarse como $c_2(X)-2+2b^1(X)$, donde
$b^1(X)=2\dim(\Pic^0_{X/k})$, y $c_2(X)$ es la segunda clase de Chern, la
característica de Euler--Poincaré y el número de intersección de la diagonal
de $X\times X$ consigo misma. Así, $b^2(X)$ tiene una interpretación
independiente de la cohomología. La desigualdad
$\rang(C^1_{\rm num}(X))\leq b^2(X)$, en esta forma y para
$\operatorname{car}(k)>0$, se debe, antes de la cohomología étale, a IGUSA,
quien utilizó resultados profundos sobre la estructura del grupo fundamental
moderado de una curva algebraica privada de algunos puntos.

\textbf{Teorema 5.4.} \emph{Sean $k$ un cuerpo algebraicamente cerrado y
$X$ un $k$-esquema propio. Existe entonces un número finito de curvas
cerradas íntegras $Y_1,\ldots,Y_p$ en $X$ con la propiedad siguiente: para
que una parte $\Lambda$ del esquema $\Pic_{X/k}$ sea cuasicompacta, es
necesario y suficiente que los números de intersección
$\langle\lambda.Y_i\rangle$ de los $\lambda\in\Lambda$ permanezcan
acotados. Además, se puede tomar como $p$ el rango del grupo de Néron--Severi
$\Pic_{X/k}(k)/\Pic^0_{X/k}(k)$.}

En efecto, como $\Pic^\tau_{X/k}$ es un subgrupo abierto cuasicompacto de
$\Pic_{X/k}$ (4.7), para que $\Lambda$ sea cuasicompacta es necesario y
suficiente que el conjunto correspondiente de puntos del espacio vectorial
$[\Pic_{X/k}(k)/\Pic^\tau_{X/k}(k)]\otimes_\Z\Q$ sea finito. Pero la
$\tau$-equivalencia coincide con la equivalencia numérica
(4.6 (i) $\Longleftrightarrow$ (iii)), por lo que entre las curvas de $X$ se
pueden elegir curvas cuyas imágenes formen una base del espacio dual; de
aquí la afirmación.

\subsection*{6. Apéndice: estudio de los $(b)$-haces sobre
$P=\mathbf P^N_k$}

\textbf{6.1.} Sean $k$ un cuerpo algebraicamente cerrado y
$P=\mathbf P^N_k$. Sean $F$ un haz coherente sobre $P$, con polinomio de
Hilbert
\[
\chi(F(n))=\sum_{i=0}^{r}a_i\binom{n+i}{i},
\]
y
\[
(c)=(c_0,\ldots,c_r)
\]
una sucesión de enteros $\geq(a)$. Sean
\[
a_{m,i}=\sum_{j=0}^{r-i}a_{j+i}\binom{m-1+j}{j}\quad\text{y}\quad
c_{m,i}=\sum_{j=0}^{r-i}c_{j+i}\binom{m-1+j}{j}.
\]
Finalmente, para todo entero $q\geq1$, denotemos el mayor subhaz coherente de
$F$ con soporte de dimensión $\leq q-1$ por $N_q$, y sea $F_q=F/N_q$.

\textbf{Lema 6.2.} \emph{Bajo las condiciones de 6.1, consideremos una
sucesión exacta}
\[
0\longrightarrow I\longrightarrow\mathcal O_P^{\oplus M}
\longrightarrow F(m)\longrightarrow0
\]
\emph{con $m\geq0$. Entonces $I$ es $p$-regular, con
$p=P_r(c_{m,0},\ldots,c_{m,r})$, donde $P_r$ es el $r$-ésimo
$(b)$-polinomio (1.9).}

En efecto, $\mathcal O_P$ es evidentemente un $(0,\ldots,0,1)$-haz
($N+1$ términos), y por tanto $I$ es un $(0,\ldots,0,M)$-haz según 1.6.
Gracias a 2.10,
$\chi(I(n))=M\binom{n+N}{N}-
\sum_{i=0}^{r}a_{m,i}\binom{n+i}{i}$. La afirmación resulta, pues, de 1.11 y
1.10.

\textbf{Proposición 6.3.} \emph{Bajo las condiciones de 6.1, supongamos que
$F$ es un $(b)$-haz. Mediante polinomios universales en los $a_i$, $b_i$ y
$N$, se pueden calcular enteros $m,M,m_1$ y $M_1$ tales que existe una sucesión
exacta}
\[
\mathcal O_P(-m_1)^{\oplus M_1}\longrightarrow
\mathcal O_P(-m)^{\oplus M}\longrightarrow F\longrightarrow0.
\]
\emph{Además, $N$ solo interviene en el cálculo de $M_1$.}

En efecto, se toma $m=P_{r-1}(c_0,\ldots,c_{r-1})$, con $c_i=b_i-a_i$, y
$M=\sum a_i\binom{m+i}{i}$, de modo que, en virtud de 1.11 y 1.3, existe
una sucesión exacta
\[
0\longrightarrow I\longrightarrow\mathcal O_P^{\oplus M}
\longrightarrow F(m)\longrightarrow0.
\]
Del mismo modo, teniendo en cuenta 6.2, se toma $m_1=m+p$, con
$p=P_r(\sisi{a_{m,0},\ldots,a_{m,r}}{c_{m,0},\ldots,c_{m,r}})$, y
$M_1=M\cdot\binom{p+N}{N}-\sum a_i\binom{m_1+i}{i}
(=\chi(I(p)))$.

\textbf{Teorema 6.4.} \emph{Bajo las condiciones de 6.1, supongamos que $F$
es un cociente de $\mathcal O_P(-m)^{\oplus M}$, con $m\geq0$. Sea
$b_i=P_{r-i}(c_{m,i},\ldots,c_{m,r})$ para $i=0,\ldots,r$, donde $P_s$ es
el $s$-ésimo $(b)$-polinomio (1.9); sea $b=b_1+m-1$ y sea
$B=\sum_{i=1}^{r}a_i\binom{b+i}{i}$. Entonces:}
\begin{enumerate}[label=(\roman*)]
\item \emph{$F$ es $b$-regular.}
\item \emph{$-B\leq a_0=h^0(F(b))-B$.}
\item \emph{$F$ es un $(b_0,\ldots,b_r)$-haz.}
\end{enumerate}

En efecto, tomemos una sección $\sigma$ de $\mathcal O_P(1)$ tal que se
obtenga un diagrama conmutativo de filas y columnas exactas:
\[
\begin{tikzcd}[column sep=large,row sep=large]
&0 \arrow[d]&0 \arrow[d]&0 \arrow[d]&\\
0 \arrow[r]&I(-1) \arrow[r] \arrow[d]&I \arrow[r] \arrow[d]&J \arrow[r]
\arrow[d]&0\\
0 \arrow[r]&\mathcal O_P(-1)^{\oplus M} \arrow[r] \arrow[d]&
\mathcal O_P^{\oplus M} \arrow[r] \arrow[d]&
\mathcal O_Y^{\oplus M} \arrow[r] \arrow[d]&0\\
0 \arrow[r]&F(m-1) \arrow[r] \arrow[d]&F(m) \arrow[r] \arrow[d]&
G(m) \arrow[r] \arrow[d]&0\\
&0&0&0&
\end{tikzcd}
\]
Como $\chi(G(n))=\sum_{i=0}^{r-1}a_{i+1}\binom{n+i}{i}$ por 1.7, $J$ es
$b_1$-regular por 6.2; por tanto, $H^q(I(p))=0$ para $q\geq2$ y
$p\geq b_1-q$ por 1.4 (i). La sucesión exacta
$H^q(\mathcal O_P^{\oplus M}(p))\longrightarrow H^q(F(m+p))
\longrightarrow H^{q+1}(I(p))$ muestra entonces que $F(m)$ es
$(b_1-1)$-regular; de aquí (i). La afirmación (i) implica que
$0\leq h^0(F(b))=\chi(F(b))=
\sum_{i=0}^{r}a_i\binom{b+i}{i}$; de aquí (ii). Según 2.10 y 1.9,
\[
\chi(F(b))=\sum_{i=0}^{r}a_{m,i}\binom{b_1-1+i}{i}
\leq b_1+\sum_{i=0}^{r}c_{m,i}\binom{b_1-1+i}{i}=b_0.
\]
Por tanto, $h^0(F(-1))\leq h^0(F(b))\leq b_0$. Por recurrencia sobre $r$,
se puede suponer que $G$ es un $(b_1,\ldots,b_r)$-haz; de aquí (iii),
gracias a 1.6.2.

\textbf{Lema 6.5.} \emph{Bajo las condiciones de 6.1, supongamos que $F$ no
tiene ningún ciclo primo asociado de dimensión $0$.}
\begin{enumerate}[label=(\roman*)]
\item \emph{Si $h^0(F)\geq1$, entonces $h^0(F(-1))\leq h^0(F)-1$.}
\item \emph{$H^0(F(-n))=0$ para $n\geq h^0(F)$.}
\item \emph{Sean una sucesión exacta
$0\to F(-1)\to F\to G\to0$ y un entero $n_0\geq0$ tales que
$H^0(G(-n_0))=0$. Entonces $H^0(F(-n_0))=0$ y
$h^0(F)\leq n_0.h^0(G)$.}
\end{enumerate}
En efecto, sea $\sigma\in H^0(F)$, $\sigma\ne0$. Entonces $\sigma$ define
un subhaz no nulo $F'=\sisi{\mathcal O_X}{\mathcal O_P}.\sigma$ de
$F$, por lo que $s=\dim\operatorname{supp}(F')$ es $\geq1$. Tomemos un diagrama
conmutativo de filas y columnas exactas:
\[
\begin{tikzcd}[column sep=large,row sep=large]
&0 \arrow[d]&0 \arrow[d]&0 \arrow[d]&\\
0 \arrow[r]&F'(-1) \arrow[r] \arrow[d]&F' \arrow[r] \arrow[d]&G'
\arrow[r] \arrow[d]&0\\
0 \arrow[r]&F(-1) \arrow[r]&F \arrow[r]&G \arrow[r]&0
\end{tikzcd}
\]
De $s\geq1$ resulta $G'\ne0$; por consiguiente, $F'$ no está contenido en
$F(-1)$, es decir, la sección $\sigma$ no pertenece a la imagen de la
inyección $H^0(F(-1))\to H^0(F)$; de aquí (i). La afirmación (ii) resulta
inmediatamente de (i), teniendo en cuenta que $h^0(F)=0$ implica
evidentemente $h^0(F(-1))=0$, pues $h^0(F(-1))\leq h^0(F)$. En (iii), se
obtiene
$0\leq h^0(F(n))-h^0(F(n-1))\leq h^0(G(n))$. Por tanto, para
$n\geq n_0$, $h^0(F(-n_0))=h^0(F(-n))$, pues $h^0(G(-n_0))=0$; de aquí
$h^0(F(-n_0))=0$ por (ii). Por consiguiente,
$h^0(F(p-n_0))\leq h^0(G(1-n_0))+\cdots+h^0(G(p-n_0))
\leq ph^0(G(p-n_0))$ para $p\geq1$; de aquí (iii).

\textbf{Proposición 6.6.} \emph{Bajo las condiciones de 6.1, supongamos que
$F$ es un $(b_0,\ldots,b_r)$-haz. Entonces $F_q$ es un
$((b_{q-1})^q,\ldots,(b_{q-1})^2,b_{q-1},b_q,\ldots,b_r)$-haz.}
En efecto, sea $\sigma$ una sección «general» de
$\sisi{\mathcal O_X(1)}{\mathcal O_P(1)}$; entonces
$\sigma$ define un diagrama conmutativo de filas y columnas exactas:
\[
\begin{tikzcd}[column sep=large,row sep=large]
&0 \arrow[d]&0 \arrow[d]&0 \arrow[d]&\\
0 \arrow[r]&N_q(-1) \arrow[r] \arrow[d]&N_q \arrow[r] \arrow[d]&G'
\arrow[r] \arrow[d]&0\\
0 \arrow[r]&F(-1) \arrow[r] \arrow[d]&F \arrow[r] \arrow[d]&G \arrow[r]
\arrow[d]&0\\
0 \arrow[r]&F_q(-1) \arrow[r] \arrow[d]&F_q \arrow[r] \arrow[d]&G''
\arrow[r] \arrow[d]&0\\
&0&0&0&
\end{tikzcd}
\]
donde $G$ es un $(b_1,\ldots,b_r)$-haz (1.6 (i)) y
$G''=G_{q-1}$ (cuestión de profundidad). Supongamos $q=1$. Entonces
$\operatorname{dim{.}supp}(N_q)=0$; por tanto, $H^1(N_q(-1))=0$ y
$h^0(\sisi{F_q(-1}{F_q(-1)})\leq h^0(F(-1))\leq b_0$. Además, $G'=0$,
de modo que $G\simeq G''$. Así, $F_1$ es efectivamente un
$(b_0,\ldots,b_r)$-haz. Además, 6.5 (ii) muestra que
$H^0(F_1(-b_0))=0$. Si $q\geq2$, se puede suponer que $G_{q-1}$ es un
$((b_{q-1})^{q-1},\ldots,b_{q-1},\ldots,b_r)$-haz y que
$H^0(G_{q-1}(-b_{q-1}))=0$, razonando por recurrencia sobre $q$. Entonces 6.5 muestra que
$H^0(F_q(-b_{q-1}))=0$ y $h^0(F_q(-1))\leq(b_{q-1})^q$; de aquí la
afirmación.

\textbf{Teorema 6.7.} \emph{Existen dos sucesiones
$\{A_i(X_0,\ldots,X_i;Y)\}$ y
$\{A_i^{(q)}(X_0,\ldots,X_q;Y)\}$ de polinomios con las propiedades
siguientes. Se suponen satisfechas las condiciones de 6.1 y que $F$ es un
cociente de $\mathcal O_P(-m)^{\oplus M}$, con $m\geq0$.}
\begin{enumerate}[label=(\roman*)]
\item \emph{Si $F$ es un $(b_0,\ldots,b_r)$-haz, entonces, para
$i=0,\ldots,r$,}
